\section{Dilute-Medium vs. Dense-Medium Formulation\label{sec:diluteVsDense}}
From the results presented in
\Cref{sec:gasComposition,ch4:speciesStudy,ch4:frozenVsNonequilibrium}, one
question that remains open is whether the dilute-medium formulation of the index of refraction
\Cref{eq:diluteRefractiveIndex} remains valid for hypersonic flows, or if there
are cases in which it breaks down and the dense-medium formulation 
\Cref{eq:denseRefractiveIndex} must be applied. \Cref{fig:diluteDenseDelta}
shows the limits of the dilute-medium formulation.
\begin{figure}[H]
    \makebox[\textwidth][c]{
        \includegraphics[width=0.9\textwidth]{chapter4/diluteVsDense/dilute_dense_delta}
    }
    \caption{Dilute-medium vs. dense-medium formulation}\label{fig:diluteDenseDelta}
\end{figure}

From \Cref{fig:diluteDenseDelta}, it appears that the dilute-medium formulation is
valid as long as $n - 1 < 10^{-2}$. For all results presented in
\Cref{ch:ch4}, the maximum value of $n - 1$ is less than $10^{-4}$,
confirming the validity of the dilute-medium formulation. However, one could consider
a hypersonic case in which the dilute-medium formulation may break down.
\Cref{fig:1kPa,fig:1MPa,fig:1GPa} shows a comparison between the dilute-medium
and dense-medium formulation for standard air at different pressures.
\begin{figure}[H]
    \makebox[\textwidth][c]{
    \includegraphics[width=0.9\textwidth]{chapter4/diluteVsDense/dilute_dense_pressure_1E+03}
    }
    \caption{Dilute-medium vs. dense-medium formulation at $1\Units{kPa}$}\label{fig:1kPa}
\end{figure}
\begin{figure}[H]
    \makebox[\textwidth][c]{
    \includegraphics[width=0.9\textwidth]{chapter4/diluteVsDense/dilute_dense_pressure_1E+06}
    }
    \caption{Dilute-medium vs. dense-medium formulation at
    $1\Units{MPa}$}\label{fig:1MPa}
\end{figure}
\begin{figure}[H]
    \makebox[\textwidth][c]{
    \includegraphics[width=0.9\textwidth]{chapter4/diluteVsDense/dilute_dense_pressure_1E+09}
    }
    \caption{Dilute-medium vs. dense-medium formulation at
    $1\Units{GPa}$}\label{fig:1GPa}
\end{figure}


Results in \Cref{fig:1kPa,fig:1MPa,fig:1GPa} were obtained using the \gls{HAOT}
package. For each case, a temperature vector and a fixed pressure were
specified to compute the index of refraction. The \textit{perfect gas} equation
\Cref{eq:equationOfState} was used to determine the densities, and the Kerl
formulation \Cref{eq:kerlFitting} was applied to estimate the polarizability.
The index of refraction was then calculated using both the dilute-medium and
dense-medium
formulations \Cref{eq:diluteRefractiveIndex,eq:denseRefractiveIndex}. All
calculations correspond to a wavelength of $633\Units{nm}$. Additional
calculations for a wavelength of $1000\Units{nm}$ are given in \crefApendex{app:diluteDenseFormulation}. \\

The results presented in \Cref{fig:1kPa,fig:1MPa,fig:1GPa} can be shown as a
function of Mach number, given in \Cref{fig:633Mach}, and altitude, given in
\Cref{fig:633AltitudeH}, for a laser with a wavelength of $633\Units{nm}$. These
plots illustrate how the flow conditions influence the optical properties,
allowing a direct comparison of the effects of compressibility and density
variations on the laser propagation at different altitudes and Mach numbers.
Additional calculations for a wavelength of $1000\Units{nm}$ are given in
\crefApendex{app:diluteDenseFormulation}.\\

\Cref{fig:633Mach} shows the pre-shock Mach number at an altitude of
$10\Units{km}$ for a laser with a wavelength of $633\Units{nm}$. As expected,
the index of refraction increases with increasing Mach number; however, the
overall change remains small, and the dilute-medium formulation continues to be
valid. This behavior indicates that, despite the compressibility effects
associated with higher Mach numbers, the flow remains sufficiently low-density
for the assumptions of the dilute approximation to hold, ensuring that
\gls{AO} distortions are minimal in this regime.
\begin{figure}[H]
    \makebox[\textwidth][c]{
    \includegraphics[width=0.9\textwidth]{chapter4/diluteVsDense/dilute_dense_mach_633.pdf}
    }
    \caption{Dilute-medium vs. dense-medium formulation as a function of
    Mach at a wavelength of $633\Units{nm}$}\label{fig:633Mach}
\end{figure}

\Cref{fig:633AltitudeH} shows the index of refraction as a function of
altitude. As expected, the index of refraction decreases at higher altitudes
because the flow density decreases with altitude. This trend reflects the
reduced number of particles available to interact with the laser, which
diminishes the optical path differences and, consequently, the \gls{AO} effects in the upper atmosphere.
\begin{figure}[H]
    \makebox[\textwidth][c]{
    \includegraphics[width=0.9\textwidth]{chapter4/diluteVsDense/dilute_dense_altitude_633.pdf}
    }
    \caption{Dilute-medium vs. dense-medium formulation as a function of
    altitude at wavelength of $633\Units{nm}$}\label{fig:633AltitudeH}
\end{figure}

Several assumptions underlie these results, including the use of the
\textit{perfect gas} equation of state, constant composition, and the Kerl
formulation for polarizability. The analysis indicates that a pressure of
$1\Units{GPa}$ would be required for the dilute-medium formulation to break down and
for the dense-medium formulation to become necessary. Such a pressure is far beyond
physically realistic values for air, even under extreme conditions such as
high-enthalpy hypersonic flows or shock-compressed regions, where densities
typically remain below $\mathcal{O}(10^2)\UnitsFrac{kg}{m^3}$. Therefore, for
hypersonic flows, the dilute-medium formulation of the index of refraction can be
considered an accurate and robust approximation.
